\documentclass[a4paper,11pt]{ltjsarticle}

% ========================================
% パッケージ
% ========================================
\usepackage{luatexja}
\usepackage{graphicx}
\usepackage{geometry}
\usepackage{xcolor}
\usepackage{hyperref}
\usepackage{listings}
\usepackage{enumitem}
\usepackage{booktabs}
\usepackage{longtable}
\usepackage{fancyhdr}

% ========================================
% ページ設定
% ========================================
\geometry{
  top=25mm,
  bottom=25mm,
  left=25mm,
  right=25mm
}

\setlength{\parindent}{0pt}
\setlength{\parskip}{0.7em}

% ========================================
% 色
% ========================================
\definecolor{uecblue}{RGB}{52,52,154}
\definecolor{codebg}{RGB}{245,245,245}
\definecolor{codecomment}{RGB}{80,120,80}

% ========================================
% ハイパーリンク
% ========================================
\hypersetup{
  colorlinks=true,
  linkcolor=uecblue,
  urlcolor=uecblue
}

% ========================================
% コード表示
% ========================================
\lstset{
  basicstyle=\ttfamily\small,
  backgroundcolor=\color{codebg},
  frame=single,
  breaklines=true,
  columns=fullflexible,
  showstringspaces=false,
  tabsize=2
}



% ========================================
% 文書
% ========================================
\begin{document}

% ========================================
% 表紙
% ========================================
\begin{titlepage}

\centering

\vspace*{35mm}

{\Large 電気通信大学体育会剣道部\par}

\vspace{12mm}

{\Huge\bfseries
ホームページ\\
運用・引継ぎマニュアル
\par}

\vspace{18mm}


\vfill

\begin{flushleft}
本資料は、電気通信大学体育会剣道部ホームページを、
今後の部員が継続して運用できるようにすることを目的とした
引継ぎ資料です。

ホームページの更新方法、部員情報の変更、
お知らせの追加、公開方法などについて説明します。
\end{flushleft}

\vfill

{\large 2026年9月 作成\par}

\end{titlepage}

% ========================================
% 目次
% ========================================
\tableofcontents
\clearpage




\subsection{主な更新作業}

通常の運用で行う作業は、主に以下のとおりです。

\begin{itemize}
  \item NEWS（お知らせ）の追加・更新
  \item 部員情報の追加・変更
  \item 学年・役職の更新
  \item 年間予定の更新
  \item Googleカレンダーによる稽古予定・大会予定の更新
  \item 大会・試合結果の掲載
  \item SNSや連絡先などの変更
  \item GitHubへの変更内容の保存
  \item Renderへのデプロイ
\end{itemize}



% ========================================
\section{ホームページの構成}
% ========================================

\subsection{フォルダ構成}

ホームページのファイルは、以下のような構成になっています。

\begin{lstlisting}
kendo_homepage/
|
|-- css/
|   `-- style.css
|
|-- images/
|   |-- main.jpg
|   |-- members-bg.png
|   `-- ...
|
|-- js/
|   `-- main.js
|
|-- .gitignore
|-- homepage.tex
|-- index.html
`-- news.html
\end{lstlisting}


\subsection{各ファイルの役割}

\begin{longtable}{p{0.27\textwidth}p{0.63\textwidth}}
\toprule
\textbf{ファイル・フォルダ} & \textbf{役割} \\
\midrule
\endhead

\texttt{index.html}
&
ホームページ本体です。
ABOUT、NEWS、SCHEDULE、MEMBERS、RESULTS、CONTACTなど、
トップページに表示される内容を記述しています。
\\

\texttt{news.html}
&
過去のお知らせを含むNEWS一覧ページです。
新しいお知らせを追加するときに編集します。
\\

\texttt{css/style.css}
&
ホームページ全体のデザインを管理しています。
文字サイズ、色、余白、スマートフォン表示、
部員カードなどの見た目を変更するときに編集します。
\\

\texttt{js/main.js}
&
ホームページの動作を管理しています。
部員カードのスライド、スマートフォン用メニュー、
ヘッダーの動作、試合結果ページ等への遷移時の
ローディング処理などが記述されています。
\\

\texttt{images/}
&
トップページやMEMBERS欄などで使用する画像を保存します。
\\

\texttt{homepage.tex}
&
本引継ぎマニュアルのLaTeXファイルです。
ホームページの仕様や運用方法を変更した場合は、
必要に応じてこの資料も更新してください。
\\

\bottomrule
\end{longtable}


\subsection{普段編集するファイル}

通常のホームページ運用では、主に

\begin{itemize}
  \item \texttt{index.html}
  \item \texttt{news.html}
\end{itemize}

を編集します。

\texttt{style.css} と \texttt{main.js} は、
ホームページのデザインや動作そのものを変更するときに使用します。

通常のお知らせ追加や部員情報の更新では、
基本的にこれらを変更する必要はありません。

\textbf{よく分からない状態で \texttt{style.css} や
\texttt{main.js} を大きく変更しないでください。}


% ========================================
\section{NEWS（お知らせ）の更新}
% ========================================

\subsection{NEWSの仕組み}

NEWSは、

\begin{enumerate}
  \item トップページのNEWS欄
  \item NEWS一覧ページ
\end{enumerate}

の2か所に表示されます。

トップページ側は \texttt{index.html}、
一覧ページ側は \texttt{news.html} で管理しています。

したがって、新しいお知らせを掲載するときは、
原則として\textbf{両方のファイルを更新します}。


\subsection{トップページにNEWSを追加する}

\texttt{index.html} を開き、

\begin{lstlisting}
<div class="news-list">
\end{lstlisting}

の中にある既存のNEWSを確認します。

1件のお知らせは、以下のような形式です。

\begin{lstlisting}[language=HTML]
<article class="news-item">
  <p class="news-date">2026.09.20</p>
  <p class="news-title">
    第75回関東学生剣道優勝大会が開催されましたので、
    結果をアップしました。
  </p>
</article>
\end{lstlisting}

新しいお知らせを追加する場合は、
既存の \texttt{article} をコピーして、
日付と文章を変更してください。


\subsection{新しいNEWSを一番上に追加する}

NEWSは新しいものから順番に表示します。

例えば、2026年10月1日のお知らせを追加する場合は、

\begin{lstlisting}[language=HTML]
<article class="news-item">
  <p class="news-date">2026.10.01</p>
  <p class="news-title">
    10月の稽古予定を更新しました。
  </p>
</article>
\end{lstlisting}

をNEWS一覧の一番上に追加します。


\subsection{NEWS一覧ページにも追加する}

次に \texttt{news.html} を開きます。

現在は、年ごとに

\begin{lstlisting}[language=HTML]
<div class="news-year">

  <h2>2026</h2>

  <div class="news-list">

    ...

  </div>

</div>
\end{lstlisting}

という形でNEWSを管理しています。

同じ年のお知らせであれば、
その年の \texttt{news-list} の一番上に追加してください。


\subsection{年が変わった場合}

新しい年になった場合は、
新しい \texttt{news-year} を作成します。

例えば2027年の場合は、

\begin{lstlisting}[language=HTML]
<div class="news-year">

  <h2>2027</h2>

  <div class="news-list">

    <!-- 2027年のお知らせをここに追加 -->

  </div>

</div>
\end{lstlisting}

を2026年のNEWSより上に追加します。


\subsection{更新時の注意}

NEWSを追加したら、必ず以下を確認してください。

\begin{itemize}
  \item 日付が正しいか
  \item 誤字・脱字がないか
  \item 新しいNEWSが一番上になっているか
  \item \texttt{index.html} と \texttt{news.html} の両方を更新したか
  \item PC表示で崩れていないか
  \item スマートフォン表示で崩れていないか
\end{itemize}


% ========================================
\section{部員情報の更新}
% ========================================

\subsection{部員情報の場所}

部員情報は \texttt{index.html} の
MEMBERSセクションで管理しています。

部員は学年ごとに分けられており、

\begin{lstlisting}[language=HTML]
<div class="grade-block">
\end{lstlisting}

が1つの学年を表します。

その中に、各部員の

\begin{lstlisting}[language=HTML]
<article class="member-card">
\end{lstlisting}

が並んでいます。


\subsection{部員カードの構造}

部員1人分は、おおよそ以下の形式です。

\begin{lstlisting}[language=HTML]
<article class="member-card">

  <p class="member-role">主将</p>

  <h4>神山 遼太郎</h4>

  <dl>

    <div>
      <dt>所属学類</dt>
      <dd>Ⅱ類</dd>
    </div>

    <div>
      <dt>段位</dt>
      <dd>三段</dd>
    </div>

    <div>
      <dt>出身高校</dt>
      <dd>宇都宮高等学校</dd>
    </div>

  </dl>

</article>
\end{lstlisting}

新しい部員を追加するときは、
既存のカードを1つコピーし、
必要な情報を書き換える方法が簡単です。


\subsection{役職がない場合}

役職がない部員についても、
レイアウトをそろえるため

\begin{lstlisting}[language=HTML]
<p class="member-role">　</p>
\end{lstlisting}

の行自体は残してください。


\subsection{年度替わり}

年度が変わった際には、MEMBERSについて以下を確認します。

\begin{enumerate}
  \item 各学年を1つ上げる
  \item 新1年生の欄を作成する
  \item 新入部員のカードを追加する
  \item 主将、主務、会計などの役職を変更する
  \item 段位に変更がある場合は更新する
  \item 卒業した部員の表示を整理する
\end{enumerate}

特に役職変更を忘れないよう注意してください。


% ========================================
\section{次に作成する項目}
% ========================================

今後、本資料には以下を追加します。

\begin{itemize}
  \item Googleカレンダーの更新方法
  \item 年間予定の更新方法
  \item Git・GitHubによる更新方法
  \item Renderへのデプロイ方法
  \item 試合結果ページの運用方法
  \item MEMBERS ONLYの運用方法
  \item 使用しているサービス・アカウント
  \item 年度替わりチェックリスト
  \item トラブル時の対処方法
  \item 引継ぎ時の注意事項
\end{itemize}


\end{document}